\documentclass[11pt,letterpaper]{article}
\usepackage[margin=1in]{geometry}
\usepackage{amsmath,amssymb,amsthm,mathrsfs,mathtools,booktabs,array,microtype}
\usepackage[hidelinks]{hyperref}
\hypersetup{pdftitle={ChannelComp Optimality and Larger Function Classes},pdfauthor={Research note prepared for Saeed Razavikia}}
\newtheorem{theorem}{Theorem}
\newtheorem{proposition}[theorem]{Proposition}
\newtheorem{corollary}[theorem]{Corollary}
\theoremstyle{remark}
\newtheorem{remark}[theorem]{Remark}
\newcommand{\enc}{\mathscr E}
\newcommand{\sq}{\mathrm{sq}}
\newcommand{\R}{\mathbb R}
\newcommand{\C}{\mathbb C}
\title{ChannelComp Optimality and Larger Function Classes}
\author{Research note prepared for Saeed Razavikia}
\date{13 September 2026}
\begin{document}
\maketitle
\begin{abstract}
We distinguish optimality for the squared-distance ChannelComp criterion from optimality for its Gaussian exponential counterpart. The previously derived quadratic and quartic formulas solve the former under a continuous-input, shared-encoder, pointwise peak constraint. We now give an explicit globally optimal encoder for the Gaussian exponential criterion on a broader infinite-dimensional target class, for every $K\ge2$ in a specified low-SNR regime. The class includes the mean multiplied by the second moment and the mean attenuated by variance. A separate theorem gives explicit squared-distance optima for polynomials of arbitrary degree in the mean and second moment, under nonnegative-coefficient and common-parity conditions. Complete certificates and a finite-alphabet counterexample clarify the scope of these claims. These are research derivations; historical novelty and independent proof review remain outstanding.
\end{abstract}

\section{Precisely which ChannelComp problem is solved?}
ChannelComp uses all-pair inequalities of the form
$$
 D\!\left(|r_\enc(\mathbf s)-r_\enc(\mathbf t)|\right)
 \ge\lambda|f(\mathbf s)-f(\mathbf t)|^2,
 \qquad r_\enc(\mathbf s)=\sum_{k=1}^K\enc(s_k).
$$
Noise-aware ChannelComp changes $D$ with the noise model \cite{noiseaware}. The supplied TCOM draft specifically uses $D(d)=e^{ad^2}$ with $a=1/(4N_0)$ for $N\sim\mathcal{CN}(0,N_0)$. Define
\begin{equation}
 \Lambda_D(\enc;f)=\inf_{\Delta f\ne0}\frac{D(|\Delta r_\enc|)}{|\Delta f|^2},
 \qquad \sup_{s\in[-1,1]}|\enc(s)|^2\le P.
 \label{eq:score}
\end{equation}

\begin{center}
\small
\begin{tabular}{@{}p{.28\textwidth}p{.62\textwidth}@{}}
\toprule
Criterion or model & Status of the preceding encoder addendum\\
\midrule
$D(d)=d^2$ & Explicit quadratic and quartic optima over continuous shared peak-admissible encoders, under the stated target assumptions.\\[1mm]
$D(d)=e^{ad^2}-1$ & Same maps and a scaled score by the local-quadratic transfer theorem.\\[1mm]
$D(d)=e^{ad^2}$ & Those formulas alone do not establish optimality. The Gaussian theorem below supplies a new explicit class and regime.\\[1mm]
Finite alphabet or a different power constraint & Optimality does not automatically transfer from the continuous peak-constrained problem.\\[1mm]
Gaussian MSE & Not solved by any of these separation claims.\\
\bottomrule
\end{tabular}
\end{center}

The finite-alphabet formulation in \cite[Eq.~(10)]{noiseaware} uses a norm constraint on the constellation vector. This differs from a pointwise peak constraint. Even with the same peak constraint, removing input values can remove the pairs that certify a continuum optimum.

For example, let $K=4$, $s_k\in\{-1,1\}$, and $V=1-(\sum_k s_k/4)^2$. The antipodal map $\enc(\pm1)=\pm\sqrt P$ has optimal binary squared score $64P/9$. Restricting the continuous variance-optimal parabola to this alphabet instead gives
$$
 \left(\frac4{\sqrt3}-\frac43\right)\frac{64P}{9}
 \simeq6.9409261755P<\frac{64P}{9}.
$$
Indeed, for any two-symbol map, received distances are proportional to $|\enc(1)-\enc(-1)|\le2\sqrt P$; the antipodal choice maximizes every such distance. This is a concrete reason to qualify ``optimal in the ChannelComp sense'' by the distance, domain, and power constraint.

\section{A larger class with an explicit Gaussian ChannelComp optimum}
Write
$$
 m(\mathbf s)=\frac1K\sum_k s_k,\qquad
 q(\mathbf s)=\frac1K\sum_k s_k^2,\qquad V=q-m^2,
$$
and let $\Omega_K$ be their attainable $(m,q)$ set. It satisfies $|m|\le1$ and $m^2\le q\le1$.

Fix any continuous $w:\Omega_K\to[0,1]$ with
$$
 w(1,1)=w(-1,1)=1.
$$
No convexity, monotonicity, or differentiability of $w$ is required. For fixed $A\in\R$ and $B\ne0$, consider
\begin{equation}
 f(\mathbf s)=A+B\,m(\mathbf s)w\bigl(m(\mathbf s),q(\mathbf s)\bigr).
 \label{eq:gaussclass}
\end{equation}
The target is fixed before encoder optimization.

\begin{theorem}[An all-$K$ Gaussian exponential certificate]
\label{thm:gaussian}
Let $K\ge2$, $P>0$, and $a>0$. For every target \eqref{eq:gaussclass}, if
\begin{equation}
 aK^2P\le\frac14,
 \label{eq:regime}
\end{equation}
then every encoder in the family
\begin{equation}
 \boxed{\enc_\beta(s)=\sqrt P\,[s+i\beta(1-s^2)]},
 \qquad0<\beta\le1,
 \label{eq:gaussianmap}
\end{equation}
is globally optimal for $D(d)=e^{ad^2}$, with
\begin{equation}
 \boxed{\Lambda_a^\star(f;P)=\frac{e^{4aK^2P}}{4B^2}}.
 \label{eq:gaussianvalue}
\end{equation}
The converse covers every peak-admissible encoder, even if it is discontinuous or fails exact computation. The proposed encoders are continuous and compute the target exactly.
\end{theorem}
\begin{proof}
\emph{Peak feasibility and exactness.}
The full peak constraint follows from
$$
 1-|\enc_\beta(s)|^2/P
 =(1-s^2)\,[1-\beta^2(1-s^2)]\ge0.
$$
The receiver recovers
$$
 m=\frac{\operatorname{Re}r}{K\sqrt P},\qquad
 q=1-\frac{\operatorname{Im}r}{K\beta\sqrt P},
$$
and evaluates \eqref{eq:gaussclass}.

\emph{Matching converse.}
The all-$1$ and all-$-1$ inputs have target gap $2|B|$. Their received distance for any peak-admissible map is at most $2K\sqrt P$. This pair upper-bounds the score by \eqref{eq:gaussianvalue}. The proposed map attains this distance.

\emph{Every other input pair satisfies the same bound.}
Normalize the output gap by $\delta=|\Delta f|/|B|$ and put $c=aK^2P\le1/4$. Let the means of the two configurations be $m,n$. Their normalized outputs $g=mw(m,q)$ and $\widetilde g=nw(n,\widetilde q)$ have the same signs as their respective means, or are zero, and $|g|\le|m|$, $|\widetilde g|\le|n|$.

If $mn\le0$, then $\delta\le|m-n|\le2$ and the real part of the received difference gives $d\ge K\sqrt P\,|m-n|\ge K\sqrt P\,\delta$. The function $e^{cx^2}/x^2$ is nonincreasing on $0<x\le2$, because $cx^2\le1$. Hence
$$
 \frac{e^{ad^2}}{|\Delta f|^2}
 \ge\frac{e^{c\delta^2}}{B^2\delta^2}
 \ge\frac{e^{4c}}{4B^2}.
$$
If $mn>0$, both normalized outputs belong either to $[0,1]$ or to $[-1,0]$. Thus $\delta\le1$, and
$$
 \frac{e^{ad^2}}{|\Delta f|^2}
 \ge\frac1{B^2}
 \ge\frac{e^{4c}}{4B^2},
$$
where the last inequality follows from $e^{4c}\le e<4$. These cases cover all pairs with nonzero target gap and prove global optimality.
\end{proof}

This class has infinitely many functional degrees of freedom. Examples, with $A=0,B=1$, include
\begin{center}
\begin{tabular}{@{}ll@{}}
\toprule
Target & Permitted parameters\\
\midrule
$m q^r$ & $r>0$\\
$m\exp[-\eta(q-m^2)]$ & $\eta\ge0$\\
$m/[1+\eta(q-m^2)]$ & $\eta\ge0$\\
$\tfrac12m[1+\cos(\omega(q-m^2))]$ & $\omega\in\R$\\
\bottomrule
\end{tabular}
\end{center}
The second and third examples attenuate the mean according to the variance. The oscillatory example illustrates that the certificate does not require a monotone readout.

Under the TCOM convention $a=1/(4N_0)$, condition \eqref{eq:regime} is $K^2P/N_0\le1$. It is a sufficient low-SNR regime and becomes more restrictive as $K$ increases at fixed $P,N_0$. We do not claim all-SNR optimality or that this regime is sharp for each individual target. Identical exponential scores across $\beta$ also do not imply identical MSE or conditioning as $\beta$ tends to zero.

\begin{corollary}[An arbitrary differentiable local feature]
\label{cor:generalfeature}
The same Gaussian optimum holds with the second moment replaced by $t=K^{-1}\sum_k\phi(s_k)$, where $\phi\in C^1([-1,1])$. Suppose $0\le w(m,t)\le1$ is continuous on the attainable moment set and $w(\pm1,\phi(\pm1))=1$. Choose any $L>0$ with $|\phi'(s)|\le L$, and let
$$
 \ell_\phi(s)=\frac{1-s}{2}\phi(-1)+\frac{1+s}{2}\phi(1).
$$
For $f=A+Bmw(m,t)$, $K\ge2$, and $aK^2P\le1/4$, a globally optimal exact encoder is
\begin{equation}
 \boxed{\enc_{\phi,\beta}(s)=\sqrt P\left[
 s+i\frac{\beta}{2L}\bigl(\phi(s)-\ell_\phi(s)\bigr)\right]},
 \quad0<\beta\le1.
 \label{eq:generalfeature}
\end{equation}
Its score is $e^{4aK^2P}/(4B^2)$.
\end{corollary}
\begin{proof}
The residual $h=\phi-\ell_\phi$ vanishes at both endpoints and has derivative bounded by $2L$. Consequently
$$
 |h(s)|\le2L\min\{1-s,1+s\}\le2L\sqrt{1-s^2},
$$
so \eqref{eq:generalfeature} satisfies the peak constraint. Its real part reveals $m$, and its imaginary part reveals
$$
 t=\ell_\phi(m)+\frac{2L}{K\beta\sqrt P}\operatorname{Im}r.
$$
The encoder uses antipodal endpoint symbols and has real received difference $K\sqrt P(m-n)$. Every step of Theorem~\ref{thm:gaussian} therefore applies unchanged.
\end{proof}

For any convex $C^1$ generator $\phi$ and $\eta\ge0$, this includes
$$
 f(\mathbf s)=m\exp\!\left[-\eta\left(\frac1K\sum_k\phi(s_k)-\phi(m)\right)\right].
$$
The Jensen gap is nonnegative, so the exponential weight lies in $[0,1]$ and equals one at consensus. Thus the Gaussian-optimal class includes arbitrary differentiable convex local generators, rather than only quadratic moments. The target and generator are fixed first; the encoder is then constructed from them.

\section{Arbitrary-degree polynomial targets under squared separation}
We can also enlarge the class with explicit optima for $D(d)=d^2$. Let
\begin{equation}
 f(\mathbf s)=F(m,q),\qquad
 F(m,q)=\sum_{(i,j)\in\mathcal I}a_{ij}m^i q^j,
 \label{eq:polyclass}
\end{equation}
where $\mathcal I$ is finite, $i,j$ are nonnegative integers, $a_{ij}\ge0$, and all present powers $i$ have the same parity $\epsilon\in\{0,1\}$. Assume the target depends on both arguments, equivalently
$$
 \alpha=F_m(1,1)=\sum i a_{ij}>0,\qquad
 \beta=F_q(1,1)=\sum j a_{ij}>0.
$$
The polynomial degree is unrestricted. This class includes products of moments and sums of interactions, extending the preceding moment-minus-mean class.

\begin{theorem}[Explicit optimum for moment polynomials]
\label{thm:polynomial}
Let $K\ge4$ and let \eqref{eq:polyclass} satisfy the stated conditions. Define
$$
 t=\frac{\alpha}{2\beta},\qquad
 b=\frac{2}{1-t^2+t\sqrt{t^2+2}},\qquad a_0=\sqrt{b(2-b)}.
$$
An optimal encoder over every continuous shared peak-admissible map is
\begin{equation}
 \boxed{\enc^\star(s)=\sqrt P\,[a_0s+i(bs^2-1)]},\qquad
 \boxed{\Lambda_{\sq}^\star(f)=\frac{K^2P}{\beta^2}b^2(b-1)}.
 \label{eq:polyopt}
\end{equation}
\end{theorem}
\begin{proof}
\emph{Step 1: Exactness forces both moments.}
Suppose two interior triples with statistics $(S,Q)$ and $(\widetilde S,\widetilde Q)$ have equal received sums. Complete both triples with one common input $x$ and $K-4$ zeros. Exactness implies the polynomial identity
$$
 F((S+x)/K,(Q+x^2)/K)
 =F((\widetilde S+x)/K,(\widetilde Q+x^2)/K)
$$
for all real $x$, since it holds on $(-1,1)$.

Put $c_{ij}=a_{ij}/K^{i+j}$ and $d=\max(i+2j)$ over nonzero terms. Every weighted degree $i+2j$ has parity $\epsilon$, so no term has degree $d-1$. If a highest-degree term has $i>0$, the coefficient of $x^{d-1}$ is
$$
 S\sum_{i+2j=d}i c_{ij},
$$
with a strictly positive multiplier. Thus $S=\widetilde S$. For a completion with $S+x\ne0$ and positive second moments, $F((S+x)/K,q)$ is strictly monotone in $q$: all nonzero derivative terms have the same sign. Hence $Q=\widetilde Q$.

Otherwise the highest-degree term is a pure second-moment power, so $d=2J$ and that term is $c_{0J}(Q+x^2)^J$. The coefficient of $x^{d-2}$ is $Jc_{0J}Q$ plus constants from terms of weighted degree $d-2$. It follows that $Q=\widetilde Q$. In this case all mean powers are even, and for each positive second moment $F$ is strictly increasing in $|m|$. Equality for all completions gives $(S+x)^2=(\widetilde S+x)^2$ on an interval and hence $S=\widetilde S$.

The continuous two-moment representation lemma in the preceding addendum now applies: if the sum of a continuous planar encoder for three inputs determines both their first and second moments, the encoder must be $c+us+vs^2$ with independent real vectors $u,v$. That lemma uses invariance of domain, local convolution, and differentiation along moment fibers; its full proof is retained in the accompanying prior addendum. The present application imposes no additional encoder smoothness assumption.

\emph{Step 2: Exact score of a quadratic curve.}
Let $G=\left(\begin{smallmatrix}A&C\\C&B\end{smallmatrix}\right)$ be the Gram matrix of $u,v$. On $|m|\le1$, $0\le q\le1$, the coefficient assumptions give $|F_m|\le\alpha$ and $|F_q|\le\beta$. At $(1,1)$ the gradient is $(\alpha,\beta)$, and at $(-1,1)$ it is $((-1)^{\epsilon+1}\alpha,(-1)^\epsilon\beta)$. Thus an all-pair line-integral bound gives
$$
 \Lambda_{\sq}(\enc;f)
 \ge\frac{K^2(AB-C^2)}{\alpha^2B+\beta^2A+2\alpha\beta|C|}.
$$
This bound is sharp. Choose distinct interior inputs approaching whichever consensus endpoint maximizes the gradient's dual norm. The two-moment map has rank two there. Arbitrarily small attainable moment displacements in the dual-gradient direction approach the displayed quotient. This proves equality.

\emph{Step 3: Peak-constrained optimization.}
Reflection of the source flips the sign of $\Delta S$ and preserves $\Delta Q$ and the absolute target gap. Averaging peak constraints at $s,-s$, projecting the translation onto the quadratic coefficient, and choosing the linear coefficient orthogonal to it preserves $A,B$ and removes $C$ without worsening the score. The three peak constraints at $0,\pm1$ then reduce the optimum to $A=P b(2-b)$ and $B=P b^2$, $1\le b<2$. The peak identity $|\enc_b(s)|^2/P=1-b^2s^2(1-s^2)\le1$ verifies the whole interval. Minimization of
$$
 \frac{\alpha^2}{b(2-b)}+\frac{\beta^2}{b^2}
$$
gives $\alpha^2b(b-1)=\beta^2(2-b)^2$, which yields \eqref{eq:polyopt}. Any nonexact competitor has zero squared score.
\end{proof}

For example, $F(m,q)=m q$ has $\alpha=\beta=1$ and therefore
$$
 \enc^\star(s)=\sqrt P\left[\frac{2\sqrt2}{3}s
 +i\left(\frac43s^2-1\right)\right],\qquad
 \Lambda_{\sq}^\star=\frac{16}{27}K^2P.
$$
This same target belongs to Theorem~\ref{thm:gaussian} with $w=q$. Its Gaussian-optimal map in \eqref{eq:regime} can instead be $\sqrt P[s+i(1-s^2)]$. The squared-optimal map has endpoint squared distance $(32/9)K^2P$, strictly below $4K^2P$. Consequently its Gaussian score is strictly below the Gaussian optimum for every $a>0$ in that regime. This explicitly shows why a squared-distance optimizer cannot be silently transferred to the exponential criterion.

More examples are $m^2q+\eta q^2+\zeta m^4$ for $\eta,\zeta\ge0$, with $\alpha=2+4\zeta$ and $\beta=1+2\eta$, and $m^3q^2+m q^4$, with $\alpha=4$, $\beta=6$. Affine output changes scale the score but preserve the map. A fixed interval homeomorphism may also calibrate all inputs before computing the moments.

\section{What is established and what remains}
The Gaussian theorem is a full global encoder optimum for the stated class, peak constraint, and regime; its proof does not depend on the preceding topological lemma. The polynomial theorem supplies an unrestricted-degree class with explicit squared-distance optima for $K\ge4$, conditional only on its listed mathematical hypotheses and the continuous representation lemma proved in the prior addendum. A seven-variable characterization for other targets remains a characterization until its optimum is actually determined.

The accompanying script executed 409 checks, all passing. It checks Gaussian all-pair inequalities and endpoint equality for several weights, peak identities and readouts for general local features, polynomial completion coefficients, squared-distance inequalities, and the finite-alphabet counterexample. The analytic all-pair proofs and converses establish the theorems; sampled checks do not replace them. No independent proof review or exhaustive novelty assessment was performed.

The next substantive gap for the Gaussian criterion is optimality outside \eqref{eq:regime} for fixed meaningful targets such as $m q$ or $m e^{-\eta V}$. Increasing the target class while retaining a small certified regime is useful, but it does not resolve that regime question. The continuous, finite-alphabet, peak-power, and constellation-energy models must also remain distinct.

\begin{thebibliography}{9}
\bibitem{noiseaware} S. Razavikia and C. Fischione, ``On Designing Modulation for Over-the-Air Computation---Part I: Noise-Aware Design,'' arXiv:2506.15950v2, 2025. The formulation comparison uses Eq.~(10) and the Gaussian development leading to Eq.~(18). \url{https://arxiv.org/html/2506.15950v2}.
\bibitem{prior} \emph{Optimal Continuous Encoders Beyond Quadratic Targets}, research addendum prepared for Saeed Razavikia, Sept.~13, 2026. Lemma~3 gives the continuous two-moment representation result; Theorem~7 gives the convex synthesis characterization. Included in the preceding delivered source package.
\end{thebibliography}
\end{document}
